\section{Related Work}

\subsection{Diffusion Language Models}
Diffusion models have achieved remarkable success in continuous domains such as image and audio generation~\citep{sohl2015deep,ddpm,song2020score,dit}.
Recently, there has been growing interest in adapting diffusion-based methods to discrete text generation. Mercury~\citep{mercury}, Seed Diffusion~\citep{seeddiffusion}, and Gemini Diffusion~\citep{geminidiffusion} are industry-developed diffusion language models that demonstrate ultra-fast inference capabilities.
In the open-source community, early methods~\citep{li2022diffusion,lovelace2023latent,gong2022diffuseq,han2023ssd,strudel2022self,chen2022analog} applies continious space diffusion directly to language,
while more recent methods~\citep{hoogeboom2021argmax,hoogeboomautoregressive,he2023diffusionbert,loudiscrete,ou2024your,sahoo2024simple} leverage discrete diffusion to achieve better scalability and performance.
In this paper, we mainly focus on the discrete version, also known as \textbf{masked diffusion language models}.
LLaDA~\citep{llada} trains an 8B-scale masked diffusion language model from scratch, demonstrating competitive performance with autoregressive counterparts.
DREAM~\citep{dream} adapts pre-trained autoregressive models into diffusion language models through continued training.
These methods offer a promising alternative to autoregressive generation by enabling parallel token sampling, but they typically generate sequences of fixed length and lack the flexibility of variable-length generation.

\subsection{Block Diffusion Language Models}
Block Diffusion Language Models bridge the gap between autoregressive and diffusion paradigms by decomposing generation into sequential blocks.
\citet{ma2025dkv,wu2025fast} use approximate KV caching to make a diffusion language model generate block-by-block.
BD3-LMs~\citep{bd3lms} trains a native block diffusion language model from scratch.
Fast-DLLM v2~\citep{fastdllmv2} and SDAR~\citep{sdar} adapt pre-trained autoregressive models into block diffusion models, enabling efficient conversion of existing AR checkpoints.
LLaDA 2.0~\citep{llada2} scales block diffusion language models up to 100B parameters, demonstrating the scalability of this approach.
These models support variable-length generation, enable KV caching across blocks, and allow parallel token sampling within each block.
However, their training suffers from low computational utilization due to the context duplication, which our method addresses.

\subsection{Data Efficiency of Diffusion Language Models}
Recent studies~\citep{prabhudesaidiffusion,ni2025diffusion,ni2025training} have revealed that diffusion language models are ``super data learners'' that can extract more information from the same amount of training data through diverse mask patterns.
Another related approach is \emph{complementary masking}~\citep{fastdllmv2}, which uses an inverted mask paired with a standard mask to enhance data utilization.
However, these methods suffer from efficiency drawbacks.
The former relies on implicit accumulation of diversity over repeated epochs with single-mask sampling,
while the latter doubles the batch size and processes each perturbation independently, failing to exploit the shared clean context.
In contrast, our method transforms this implicit benefit into an explicit, efficient training strategy.
By processing multiple mask patterns simultaneously with \textbf{context sharing}, we actively enforce diverse supervision within each step while eliminating the computational waste of encoding the same context repeatedly.
